\documentclass[UTF8,12pt]{ctexart}
\usepackage[a4paper,margin=24mm]{geometry}
\usepackage{amsmath,amssymb,bm,graphicx,adjustbox}
\newcommand{\fitmath}[1]{\begin{adjustbox}{max width=\linewidth}$\displaystyle #1$\end{adjustbox}}
\setlength{\parindent}{0pt}
\setlength{\parskip}{.7em}
\sloppy
\allowdisplaybreaks
\begin{document}
\section*{2000 年考研数学三 · 参考答案}
\subsection*{2000 年真题参考答案}

\subsection*{一、填空题}

（1）\(\displaystyle yf'_1+\dfrac{\displaystyle 1}{\displaystyle y} f'_2-\dfrac{\displaystyle y}{\displaystyle x^2}g'\)。（2）\(\displaystyle \dfrac{\displaystyle \pi}{\displaystyle 4e}\)。（3）\(\displaystyle 24\)。（4）\(\displaystyle [1,\allowbreak 3]\)。（5）\(\displaystyle \dfrac{\displaystyle 8}{\displaystyle 9}\)。


\subsection*{二、选择题}

（1）D。（2）B。（3）C。（4）A。（5）C。


三、\(\displaystyle y=\dfrac{\displaystyle 3}{\displaystyle 4}+\dfrac{\displaystyle 1}{\displaystyle 4}(1+2x)e^{2x}\)。


四、\(\displaystyle a^2\left(\dfrac{\displaystyle \pi^2}{\displaystyle 16}-\dfrac{\displaystyle 1}{\displaystyle 2}\right)\)。


五、（1）当 \(\displaystyle p_1=10,\allowbreak Q_1=4,\allowbreak p_2=7,\allowbreak Q_2=5\) 时，利润最大，为 52 万元。（2）当 \(\displaystyle Q_1=5,\allowbreak Q_2=4,\allowbreak p_1=p_2=8\) 时，利润最大，为 49 万元。


六、单调递增区间为 \(\displaystyle ( -\infty,\allowbreak -1 )\)，\(\displaystyle (0,\allowbreak +\infty)\)，单调递减区间为 \(\displaystyle (-1,\allowbreak 0)\)。极大值为 \(\displaystyle f(-1)=-2e^{\pi/4}\)，极小值为 \(\displaystyle f(0)=-e^{\pi/2}\)。曲线有两条渐近线：\(\displaystyle y_1=e^\pi(x-2)\)，\(\displaystyle y_2=x-2\)。


七、\(\displaystyle \sum_{n=0}^{\infty}I_n=\ln(2+\sqrt2)\)。


八、证明略。（考虑函数 \(\displaystyle F(x)=\displaystyle \int_0^x f(t)\,\mathrm dt\)，\(\displaystyle 0\le x\le\pi\)。\(\displaystyle F(0)=F(\pi)=0\)，再证明存在 \(\displaystyle (0,\allowbreak \pi)\) 内一点 \(\displaystyle \xi\)，满足 \(\displaystyle F(\xi)=0\)，对 \(\displaystyle F(x)\) 使用罗尔定理。）


九、（1）当 \(\displaystyle a\ne-4\) 时，\(\displaystyle \beta\) 可由 \(\displaystyle \alpha_1,\allowbreak \alpha_2,\allowbreak \alpha_3\) 唯一地线性表示。（2）当 \(\displaystyle a=-4\)，且 \(\displaystyle 3b-c\ne1\) 时，\(\displaystyle \beta\) 不能由 \(\displaystyle \alpha_1,\allowbreak \alpha_2,\allowbreak \alpha_3\) 线性表示。（3）当 \(\displaystyle a=-4\)，且 \(\displaystyle 3b-c=1\) 时，\(\displaystyle \beta\) 可由 \(\displaystyle \alpha_1,\allowbreak \alpha_2,\allowbreak \alpha_3\) 线性表示，但表示不唯一。\(\displaystyle \beta=k\alpha_1-(2k+b+1)\alpha_2+(2b+1)\alpha_3\)，其中 \(\displaystyle k\) 为任意常数。


十、当 \(\displaystyle a_1a_2\cdots a_n\ne(-1)^n\) 时，二次型 \(\displaystyle f(x_1,\allowbreak x_2,\allowbreak \ldots,\allowbreak x_n)\) 为正定二次型。


十一、（1）\(\displaystyle E(X)=e^{\mu+\dfrac{\displaystyle 1}{\displaystyle 2}}\)。（2）\(\displaystyle (-0.98,\allowbreak 0.98)\)。（3）\(\displaystyle (e^{-0.48},\allowbreak e^{1.48})\)。


十二、证明略。（分别计算 \(\displaystyle E(X),\allowbreak E(Y),\allowbreak E(XY)\) 及 \(\displaystyle \operatorname{Cov}(X,\allowbreak Y)\)。）

\end{document}
